\MajorSection{القسم الأول}

\Couplet{١}{أزفت الساعة؛ وخرجت الملكة الآخذة في النقصان لتملك ما بقي من الليل؛}{تتوجها لمعة نجم، ويقوم عرشها على قرص من نور رمادي كالرّماد.}{}

\Couplet{٢}{يمسح «ذيل الذئب» المشرق الآخذ في الشحوب، فيخلّف وراءه ظلمة أشد؛}{ويرفع الفجر رأسه المضيء، متنهدًا بما يشبه هبّة ريح.}{\BurtonNote{حاشية برتون: الفجر الكاذب.}}

\Couplet{٣}{تلتقط الأعالي ذلك اللمع الشرقي، فيما لا تزال المنخفضات مضطجعة في أرجوانها؛}{وتتصاعد الضبابات اللؤلؤية، زينة الصباح، كالبخور لتحيّي السماء.}{}

\Couplet{٤}{تصهل الخيل، وتئن الإبل، وتلمع المشاعل، وتتوهج مواقدها المحمولة؛}{وتتهاوى مدينة الخيام، ويقتحم الإنسان الهواء بضجيجه وضرباته.}{}

\Couplet{٥}{تنفتح الأبواب الذهبية يمينًا ويسارًا؛ وتوثب الشمس بجبين ملتهب؛}{وتذوب سحابة الندى في دفق من الضوء؛ وتغتسل الأرض السمراء بوهج الصباح.}{}

\Couplet{٦}{يسيرون ببطء ملتفّين في الفلاة، فيما يعلو نشيد النهار الفتي؛}{ويقع حزينًا على أذني المتلهفة رنين أجراس الإبل.}{}

\Couplet{٧}{عبر قفار محرقة وبرارٍ متجمدة، وعبر جبل موحش ووهدة مظلمة؛}{حيث تسكن الوحوش المرعبة والغول، وحيث يأوي أناس أشد وحشية وهولًا.}{\BurtonNote{حاشية برتون: الغول شيطان الصحراء.}}

\Couplet{٨}{ومع فرحة عابرة بالنخيل الذي يعلو ويتمايل فوق السهل المضطرم؛}{حاملًا خواطر الظل الحفيف، والنبع المتفجر، والمطر المنهمر.}{}

\Couplet{٩}{ومع عزاء قصير تمنحه الحافة الجبلية التي تداعبها أنسام لطيفة؛}{فمن رأسها المهوّى وسفحها المشجّر تطل على بحار من خضرة باردة شاحبة.}{}

\Couplet{١٠}{لهم أن يمضوا في فرح ورجاء؛ فلن تكف أرواحهم عن الاهتزاز والامتلاء}{بأحلام مسقط الرأس والقبر، وبصور جبل الله المقدس.}{\BurtonNote{حاشية برتون: عرفات قرب مكة.}}

\Couplet{١١}{أما نحن؟ فمشهد آخر يتبدّل، ووخزة أخرى تمزّق القلب؛}{لماذا نلتقي على جسر الزمان لنتبادل تحية واحدة ثم نفترق؟}{}

\Couplet{١٢}{نلتقي لنفترق؛ ولكن روحي تسأل: أنفترق لنلتقي؟ آه، أحقًا يكون ذلك؟}{العلم المحيط الذي صنعه خيال الإنسان يعرف؛ أما صانع العلم المحيط فلا يستطيع أن يعرف شيئًا.}{}

\Couplet{١٣}{لماذا لا بد أن نلتقي، ولا بد أن نفترق، ولا بد أن نحمل نير هذا «لا بد»؛}{يفرضه القدر الطاغية على ضحيته، من غير أن يستأذننا أو نأذن له؟}{}

\Couplet{١٤}{ذلك المساء البهي، المضيء، المسرور؛ وهذا الصباح الكدر، الحزين، الرمادي؛}{عجيب أن يكتب مسجّل الحياة عن هذا يومًا، وعن ذاك يومًا!}{}

\Couplet{١٥}{عيناي ودماغي وقلبي حزينة؛ وحزين حتى صميم كياني؛}{كل شيء يتعب ويتغير ويمضي وينتهي؛ آه من أذى يوم الميلاد!}{}

\Couplet{١٦}{يا أصدقاء شبابي، هذا وداع أخير! لعلنا نلتقي يومًا من الأيام؛}{لكن الرجال أنفسهم لن يلتقوا أبدًا؛ فالسنون ستجعل منا رجالًا آخرين.}{}

\Couplet{١٧}{بلغ ضوء الصباح الظهيرة، وشحب مع المساء؛ والآن الوداع!}{امضوا، وغيبوا عن حياتي كما يموت رنين جرس الجمل.}{}
